\documentclass[12pt]{extarticle}
\def\activityid{F02-M-v2}\usepackage[letterpaper,margin=.65in,top=.60in,bottom=.65in]{geometry}
\usepackage[T1]{fontenc}\usepackage[default]{sourcesanspro}
\usepackage{microtype,tikz,array,tabularx,fancyhdr,amsmath,amssymb}
\usepackage[hidelinks]{hyperref}\usetikzlibrary{calc,arrows.meta}
\setlength{\parindent}{0pt}\setlength{\parskip}{7pt}\setlength{\headheight}{0pt}\setlength{\footskip}{26pt}
\setlength{\emergencystretch}{2em}\pagestyle{fancy}\fancyhf{}\renewcommand{\headrulewidth}{0pt}\renewcommand{\footrulewidth}{.4pt}
\fancyfoot[L]{\scriptsize Bellingham Math Circle / Week 2 / \activityid}\fancyfoot[R]{\scriptsize \thepage}
\definecolor{ink}{gray}{.12}\definecolor{soft}{gray}{.95}\color{ink}
\newcommand{\PageTitle}[2]{{\fontsize{10}{12}\selectfont\bfseries WEEK 2\quad TOGGLE LAB\hfill #1\par}\vspace{3pt}{\fontsize{26}{29}\selectfont\bfseries #2\par}\vspace{2pt}}
\newcommand{\NameLine}{{\small Name: \rule{2.65in}{.4pt}\hfill Date: \rule{1.35in}{.4pt}\par}\vspace{4pt}}
\newcommand{\Task}[2]{\par\vspace{3pt}{\large\bfseries #1}\par #2\par}
\newcommand{\Subhead}[1]{\par\vspace{3pt}{\large\bfseries #1}\par}
\newcommand{\RuleBox}[1]{\par\begingroup\setlength{\fboxsep}{8pt}\colorbox{soft}{\parbox{\dimexpr\linewidth-16pt\relax}{#1}}\endgroup\par}
\newcommand{\WriteBox}[2]{\par\textbf{#1}\par\vspace{2pt}\begin{tikzpicture}\draw[black!40,line width=.5pt] (0,0) rectangle (\dimexpr\linewidth-.5pt\relax,#2);\end{tikzpicture}\par}
\newcommand{\StudentFooter}{\vfill{\small Try it with objects. Explain by pointing, talking, or drawing.}\par}
\newcommand{\LampRing}[3]{\begin{tikzpicture}[x=1in,y=1in]
\foreach \i in {1,...,#1}{\coordinate (v\i) at ({90-360*(\i-1)/#1}:#2);}
\foreach \i in {1,...,#1}{\pgfmathtruncatemacro{\j}{mod(\i,#1)+1}\draw[thick] (v\i)--node[midway,fill=white,inner sep=2pt,font=\small]{\i\j}(v\j);}
\foreach \i in {1,...,#1}{\node[circle,draw,fill=white,minimum size=.52in,inner sep=0pt,font=\bfseries] at(v\i){\i};}
\foreach \i in {#3}{\node[circle,draw,fill=ink,text=white,minimum size=.52in,inner sep=0pt,font=\bfseries] at(v\i){\i};}
\end{tikzpicture}}
\newcommand{\Lights}[1]{\begin{tikzpicture}[x=.55in,y=.55in,baseline=-.10in]\foreach \v [count=\i] in {#1}{\ifnum\v=1\fill (\i,0) circle(.16in);\else\draw (\i,0) circle(.16in);\fi}\end{tikzpicture}}
\newcommand{\ClockRing}[2]{\begin{tikzpicture}[x=1in,y=1in]\draw[black!30] (0,0) circle(#2);\pgfmathtruncatemacro{\n}{#1-1}\foreach\i in {0,...,\n}{\node[circle,draw,fill=white,minimum size=.35in,inner sep=0pt] at({90-360*\i/#1}:#2){\i};}\draw[-{Stealth},thick] (-.35,.15) arc(155:25:.38);\node[font=\small] at(0,-.18){clockwise};\end{tikzpicture}}
\newcommand{\Key}[2]{\begin{center}\renewcommand{\arraystretch}{1.55}\begin{tabular}{c|*{#1}{c}}in&#2\end{tabular}\end{center}}

\begin{document}

\PageTitle{GRADES 2--3 / 1}{Which pictures are possible?}\NameLine
\RuleBox{Start all lamps OFF. Press a line to flip \textbf{both} end lamps. Dark means ON. A lamp already ON turns OFF. Write 12 for the line joining lamps 1 and 2.}
\begin{center}\LampRing{4}{.85}{}\end{center}
\Task{1. Three targets.}{Reset to all OFF before each target. Find moves, or give a reason the target cannot happen.}
\begin{center}\renewcommand{\arraystretch}{2.1}\begin{tabular}{l|p{3.4in}}ON lamps&My moves or my reason\\\hline 1 and 2&\\1 and 3&\\1 only&\\\end{tabular}\end{center}
\Task{2. Look at one move closely.}{The two end lamps might both be OFF, both ON, or one of each. How does the number of ON lamps change in these three cases?}
\WriteBox{A rule that survives every move:}{.7in}
\StudentFooter
\newpage
\PageTitle{GRADES 2--3 / 2}{Find the whole world}
\Task{3. Make a complete collection.}{Draw every reachable picture in the boxes: a dot for each lamp, filled if ON. Keep the order 1, 2, 3, 4. The all-OFF picture counts. There may be more boxes than you need.}
\begin{center}\begin{tikzpicture}[x=1in,y=1in]\foreach \r in {0,1,2}{\foreach \c in {0,1,2,3}{\draw[black!40] (1.65*\c,-.78*\r) rectangle +(1.50,.62);}}\end{tikzpicture}\end{center}
\textbf{Our count:}\ \rule{.7in}{.4pt}\qquad How do you know none are missing?
\Task{4. Make two far-apart lamps.}{Choose any two lamps. Trace a path between them. Press each line on your path once. What happens to the lamps in the middle?}
\WriteBox{Explain why a path changes only its two ends.}{1in}
\textbf{Go further:} What if the ring has 5 lamps? Can every picture with an even number of ON lamps still be made? Pair its ON lamps and try paths.
\StudentFooter

\end{document}
